\documentclass[11pt]{article}
\usepackage[margin=1in]{geometry}
% Compiled locally with Tectonic (XeTeX). For arXiv submission, replace the
% hangul font with UnBatang (in arXiv's TeX Live via unfonts-core) or submit
% the compiled PDF.
\usepackage{fontspec}
\usepackage{kotex} % Korean text support (xetexko under XeTeX)
\IfFontExistsTF{Apple SD Gothic Neo}{\setmainhangulfont{Apple SD Gothic Neo}}{}
\usepackage{booktabs}
\usepackage{graphicx}
\usepackage{natbib}
\usepackage{url}
\usepackage{xcolor}
\usepackage{multirow}
\usepackage{caption}
\captionsetup{font=small,labelfont=bf}
\usepackage[hidelinks]{hyperref}

% NOTE TO AUTHOR: verify the romanization of your name before submission.
\title{\textbf{MunBench: A Generative Benchmark for Korean Emotional
Intelligence, Creative Writing, and Cultural Nuance in Large Language Models}}
\author{Taewook Ha\\
\texttt{hataewook12@gmail.com}\\
\small\url{https://github.com/ghandhitechnology/munbench}}
\date{July 2026}

\newcommand{\ko}[1]{#1}

\begin{document}
\maketitle

\begin{abstract}
Korean-language evaluation of large language models (LLMs) is dominated by
multiple-choice knowledge tests. These are cheap to grade and resistant to
ambiguity, but they cannot answer the question practitioners actually ask:
\emph{can this model write well in Korean?} Prior work shows that factual
knowledge of a culture correlates only weakly with the ability to apply that
culture in open-ended generation, so multiple-choice scores are a poor proxy.
We present \textbf{MunBench} (\ko{문벤치}), a fully generative benchmark that
evaluates three abilities at once: (1)~\emph{emotional intelligence} in
Korean interpersonal context, tested through multi-turn roleplay scenarios
built around \ko{눈치} (reading the unstated), \ko{정} (the deep
affective bond formed over shared time, with its attendant obligations),
and \ko{체면} (face); (2)~\emph{creative writing}, tested
through 48 commissions that each target a documented LLM weakness such as
melodrama and clich\'e; and (3)~\emph{cultural and linguistic nuance},
tested through honorific-register switching, idiom deployment, sarcasm
resolution, and paired prompts that reveal whether a model silently defaults
to Western norms when a Korean setting is implied but not stated. All 158
items were authored natively in Korean. Scoring combines an ensemble
rubric, anchored pairwise Elo, and judge-free mechanical metrics that catch
failure modes LLM judges are known to miss, including code-switching into
English. In an initial 13-model study, three results stand out. The
creative-writing track separates models that other tracks cannot: writing
scores span 1.7--8.2, while eight of thirteen models sit within one point
of each other on the EQ track. The culture-pair probe finds four models
whose scores \emph{drop} when a Korean setting is made explicit --- by
2.8 rubric points in the worst case --- a pattern no culture-specified
benchmark can see. And a judge-free tripwire catches a frontier model
code-switching out of Korean in 42\% of nuance-track outputs, a failure
its rubric scores alone would blur. Benchmark, harness,
and results are public.
\end{abstract}

\section{Introduction}

A user who asks an LLM for help in Korean is rarely asking it to answer exam
questions. They are asking it to write: a reply to a delicate message, a
eulogy, a cover letter, a story. Whether the model does this \emph{well} in
Korean --- with the right honorific register, without translated-English
phrasing, reading what the other person meant rather than what they said ---
is the ability that matters in practice, and it is almost entirely
unmeasured.

Korean LLM benchmarks overwhelmingly use multiple-choice formats: KMMLU
\citep{son2024kmmlu}, HAE-RAE Bench \citep{son2023haerae}, CLIcK
\citep{kim2024click}, KoBEST \citep{jang2022kobest}, and KorNAT
\citep{lee2024kornat} all grade selection, not production. This is a
reasonable design choice for knowledge, but it leaves generation quality
unmeasured, and there is direct evidence that the two dissociate:
Nunchi-Bench \citep{nunchibench2025} found that accuracy on factual
questions about Korean cultural practices correlates only weakly with
appropriate behavior in open-ended scenarios that depend on those same
practices. Knowing \emph{about} \ko{눈치} is not the same as having it.
The few Korean benchmarks that do score generation have used a single
generic quality scale; LogicKor, the most widely used, was retired after
frontier models saturated its 0--10 judge rubric. Meanwhile, the
English-language evaluation community has developed strong methodology for
exactly this problem --- EQ-Bench~3's judged roleplay with paired
rubric-and-Elo scoring \citep{paech2023eqbench}, bias controls for LLM
judges \citep{zheng2023judging}, and reward-model findings from LitBench
\citep{litbench2025} --- but none of it has been applied to Korean.

MunBench fills this gap. It is, to our knowledge, the first benchmark that
evaluates Korean \emph{generation} across emotional intelligence, creative
writing, and cultural nuance simultaneously, with methodology adapted from
the strongest English-language precedents and grounded in
Korean-specific failure modes. Our contributions:

\begin{enumerate}
\item \textbf{A 158-item generative benchmark}, authored natively in Korean
(zero translated items), spanning three tracks: 50 multi-turn EQ roleplay
scenarios, 48 creative-writing commissions, and 60 nuance tasks, of which 33
carry \emph{contrast variants} that isolate specific abilities
(Section~\ref{sec:design}).
\item \textbf{A three-layer scoring pipeline} --- ensemble rubric scoring,
anchored pairwise Elo, and judge-free mechanical metrics --- with position,
length, and self-preference bias controls, and an explicit code-switching
penalty motivated by evidence that LLM judges under-detect language mixing
in Korean text \citep{son2024kudge} (Section~\ref{sec:scoring}).
\item \textbf{An initial 13-model study} showing that (a) creative writing
is the discriminating track; (b) four models score \emph{higher} when
Korean context is implicit than when it is explicit, a pattern consistent
with explicit Korean framing activating degraded generation; and (c) judge-free metrics
catch large failure modes (up to 42\% code-switching) that rubric scores
alone would blur (Section~\ref{sec:results}).
\end{enumerate}

\section{Related Work}

\paragraph{Korean benchmarks.} KMMLU \citep{son2024kmmlu} builds
multiple-choice items from original Korean exams rather than translating
MMLU; HAE-RAE Bench \citep{son2023haerae} probes knowledge that
English-centric models lack; CLIcK \citep{kim2024click} covers cultural and
linguistic knowledge from Korean textbooks; KoBALT \citep{kobalt2025}
covers 24 linguistic phenomena; KorNAT \citep{lee2024kornat} measures
alignment with Korean population survey distributions. All are
selection-based. KoBBQ \citep{jin2023kobbq} and related safety datasets
grade classification. Generation-oriented efforts are rare: LogicKor scored
free-form Korean with a generic judge rubric (now saturated), and KUDGE
\citep{son2024kudge} --- a meta-benchmark --- showed that LLM-judge
reliability on Korean text tracks the judge's English capability and that
judges systematically miss factual errors and language mixing, a finding
that directly shaped our judge-free metric layer. KoCoSa
\citep{kim2024kocosa} showed context-dependent sarcasm is difficult in
Korean; our subtext tasks adopt this framing. BLEnD \citep{myung2024blend}
demonstrated the value of separating cultural awareness from language
competence.

\paragraph{Emotional intelligence and creative writing.} EQ-Bench
\citep{paech2023eqbench} evolved from emotion-intensity prediction to
judged multi-turn roleplay (version~3) with paired absolute-rubric and
pairwise-Elo scoring; we adopt this two-scale design. EmoBench
\citep{sabour2024emobench} and ToMBench \citep{chen2024tombench} test
emotional and theory-of-mind reasoning via multiple choice, and ToMBench's
native-language-first construction (author in the target language, then
translate for documentation) is the order we follow. For creative writing,
LitBench \citep{litbench2025} showed that off-the-shelf LLM judges lag
trained reward models at predicting human preferences and that
chain-of-thought can \emph{hurt} creative judging; NoveltyBench
\citep{noveltybench2025} showed aligned models lose output diversity;
WritingBench \citep{writingbench2025} used per-instance criteria. Judged
evaluation at scale rests on bias controls established by MT-Bench and
Chatbot Arena \citep{zheng2023judging,chiang2024arena}: we use both-order
judging, equal-length truncation, and anonymization throughout.

\section{Benchmark Design}
\label{sec:design}

MunBench contains 158 items in three tracks. Every item was written
natively in Korean --- none are translations --- following the
native-first construction argued for by ToMBench and by the authors of
KMMLU and CLIcK, who document systematic artifacts in translated
benchmarks. Every item also carries \ko{채점 유의사항} (judge notes): a
per-item description of what strong and weak responses look like, which is
provided to judges at scoring time.

\subsection{Track 1: \ko{감정} --- emotional intelligence (50 scenarios)}

Each scenario places the model in a role within an emotionally charged,
specifically Korean situation, spanning five domains: workplace, family
(including in-laws), friends and romance, neighbors and service
encounters, and school hierarchies. The interlocutor's three turns are
\emph{pre-scripted} and written to remain coherent after any plausible
reply, so every model faces an identical conversation --- adopting the
fixed-turn design of EQ-Bench~3 \citep{paech2023eqbench}. After the
roleplay, the model answers an out-of-character analysis question about the
interlocutor's actual mental state. Each item is tagged with the phenomena
it stresses: \ko{눈치}, \ko{정}, \ko{체면}, honorific-register pressure,
sarcasm (\ko{반어}), perceived slight (\ko{서운함}), and power asymmetry
(\ko{갑을관계}).

A representative scenario: the model plays a junior employee whose team
held a company dinner without the team leader. The leader's scripted turns
deny being hurt (``\ko{서운하다는 건 아니고}...''), invoke etiquette,
then insist on buying coffee. Strong responses read the hurt beneath the
denial without naming it to the leader's face, preserve face, and signal
future inclusion; the judge notes list characteristic failures ---
excessive apology, therapist-style emotion labeling
(``\ko{서운하셨군요}...''), and jumping to a rescheduled dinner as a
transactional fix.

\subsection{Track 2: \ko{문학} --- creative writing (48 commissions)}

Forty-eight writing commissions in seven forms: short-fiction scene,
personal essay, poetry, dialogue-driven scene, genre fiction, constraint
stories with mandatory unrelated elements in the style of
\citealp{lechmazur2024}, and one dramatic monologue --- mostly
600--1{,}200 Korean characters. Each
commission deliberately targets a documented LLM weakness: writing in a
twelve-year-old's authentic voice rather than an adult's imitation of one;
conveying loss strictly through objects and absence with direct emotion
statements banned; resisting \ko{신파} (melodrama). This
weakness-targeting philosophy follows EQ-Bench Creative Writing v3, whose
prompts are chosen to expose failure rather than to elicit typical output.

\subsection{Track 3: \ko{결} --- nuance and culture (60 tasks)}

Four task families, 15 items each, all generative:

\begin{itemize}
\item \textbf{Register switching}: rewrite or continue a scene under a
changed relationship (\ko{반말} to formal \ko{존댓말}, deferential service
register, and unscripted cases such as a same-age coworker just met).
Honorific correctness is scored as its own criterion, never folded into
fluency, following RULER/VERSE \citep{rulerverse2025}.
\item \textbf{Idiom deployment}: weave a given proverb or
four-character idiom (\ko{사자성어}) naturally into a scene ---
explaining the idiom or announcing it (``\ko{이런 말 있잖아}'') is a
failure. Trap items use near-miss idioms that models commonly confuse.
\item \textbf{Culture pairs}: the same task issued twice --- once stating
the setting is Korea, once leaving it unstated while every name and
circumstance remains Korean. The probe, adapted from Nunchi-Bench's
specified-vs.-neutral design \citep{nunchibench2025}, detects models that
default to Western norms (hugging, first names, Western funeral etiquette)
when Korean-ness is implicit.
\item \textbf{Sarcasm and subtext}: respond to a dialogue whose final
utterance means the opposite of its surface (``\ko{아니야 괜찮아, 일이
우선이지 뭐}''), in the context-dependent style of KoCoSa
\citep{kim2024kocosa}. Several items carry a \emph{sincere-version
control}: the same words in a context where they are genuinely meant, so
that a model that hallucinates hurt feelings everywhere also fails.
Over-reading subtext is an emotional-intelligence failure, not a success.
\end{itemize}

In total, 33 items carry contrast variants: all 15 culture pairs
(Korea stated vs.\ unstated), 15 idiom items (a de-contextualized control
prompt), and 3 subtext items (the sincere-version control). Only the
culture-pair variants feed the specified-vs.-implied delta of
Section~\ref{sec:results}; the other variants probe different constructs
and are scored separately. A full run produces exactly 191 outputs per
model.

\subsection{Item construction and quality control}
\label{sec:construction}

Items were drafted by a frontier LLM (Claude Opus~4.8) under detailed
track-specific style guides, in batches by domain or form. Every batch
then passed an adversarial critique stage: a separate model reviewed each
batch against six criteria (translated-sounding Korean; whether the target
phenomenon is actually embedded rather than merely named; discriminative
power; clich\'e; robustness of pre-scripted turns to arbitrary replies;
specificity of judge notes), and flagged items were revised by the
authoring model. The critique stage requested revisions in 14 of 15
batches. We are explicit that this is \emph{LLM-authored, LLM-audited}
data --- a deliberate trade-off that bought native-quality volume at low
cost, with risks discussed in Section~\ref{sec:limitations}.

\section{Scoring}
\label{sec:scoring}

Every output is scored three independent ways; each layer covers a known
blind spot of the others.

\subsection{Ensemble rubric (interpretable)}

Each track has a weighted rubric of six to seven Korean-specific criteria
--- for Track~1: reading \ko{눈치}/\ko{정}/\ko{체면} (weight 0.22),
emotional granularity (0.18), honorific correctness (0.17), restraint and
anti-sycophancy (0.16), character consistency (0.15), and situational
navigation (0.12). A generic single quality score is deliberately avoided:
it is what saturated LogicKor. An ensemble of three judges from different
model families scores every criterion 0--10; per-judge scores are stored
separately so that judge disagreement and self-preference are inspectable,
and each judge can score repeatedly (configurable iterations) to estimate
repeatability, following the reporting practice of EQ-Bench~3. Criterion
names returned by judges are matched with normalization and fuzzy
fallback, and partially parseable judgments receive renormalized partial
credit rather than being zeroed. An anti-sycophancy criterion appears in
every rubric because EQ-flavored benchmarks are especially vulnerable to
reward-hacking through excessive warmth \citep{paech2023eqbench}.

\subsection{Anchored pairwise Elo (discriminative)}

Absolute rubric scores compress as models improve. The second layer
therefore asks judges only to choose between two \emph{anonymized} outputs
for the same item, truncated to equal length (removing length bias), with
every pair judged in both label orders and averaged (removing position
bias) \citep{zheng2023judging}. A comparison is discarded --- not
averaged from one order --- if either order fails, and comparisons in
which the judge's own model appears as a contestant are excluded from the
fit entirely: anonymization hides names, not style, and a judge voting on
its own matchup would inject self-preference directly into the ratings.
Wins are fitted to Elo ratings with three \emph{anchor} models --- chosen
to span the ability range, one of them Korean-specialized --- all pinned
at a common reference rating of 1200. Pinning keeps the scale stable as
new models are added; because all anchors share one pinned value, an
anchor's Elo is a reference point rather than a measurement of that
anchor (anchors' abilities are read from their rubric scores), and
non-anchor ratings are interpreted relative to the anchor set as a whole.
Assigning distinct pinned values per anchor is planned for the full-scale
run.

\subsection{Judge-free mechanical metrics (tripwires)}

Three metrics involve no LLM opinion: (1)~a curated 184-phrase Korean
``slop'' list of LLM-overused phrasing and translationese
(\ko{마음 한켠}, \ko{심장이 쿵 내려앉}--, \ko{$\sim$에 다름 아니다}),
reported as hits per 1{,}000 characters; (2)~$n$-gram repetition and
length-specification compliance (supporting both character-count and
line/stanza specifications for poetry); and (3)~a \emph{code-switching
flag} when Latin characters exceed 2\% of the output's alphabetic
characters. The last is
motivated directly by KUDGE's finding that LLM judges under-detect
language mixing in Korean text \citep{son2024kudge}; our results confirm
the metric does work an LLM judge layer did not (Section~\ref{sec:results}).

\section{Initial Study}
\label{sec:results}

\subsection{Setup}

We evaluated 13 models (Table~\ref{tab:main}) on a stratified 47-item
subset (10 EQ scenarios covering all five domains, 13 writing commissions
covering all six forms, 24 nuance tasks covering all four families),
yielding 60 records per model including contrast variants --- 780 scored
samples total, with zero generation failures and zero rubric-judging
failures. Reasoning-capable models ran at the provider's highest
documented reasoning-effort setting.\footnote{Two models were requested
at ``max''; the serving layer maps this to \texttt{high}, its top
documented level, so all reasoning models ran at effectively the same
setting.} Judges were Claude Opus~4.6, Gemini 3.6 Flash, and GPT-5.6
Luna, one rubric iteration each. The two API-routed judges performed the
pairwise pass: 420 pairs $\times$ 2 judges $=$ 840 verdicts, each
verdict averaging two label orders (1{,}680 judge calls). Of these
verdicts, 135 failed verdict parsing and 134 involved the judge's own
model as a contestant; both classes were excluded, leaving 571
fully-order-averaged games in the Elo fit. Anchors were Claude Opus~4.8, Solar Pro~3 (the Korean-specialized
mid anchor), and GPT-5.6 Luna. Total cost was approximately \$41 in API
credits plus 780 judge calls routed through a consumer Claude
subscription. The subset design, single iteration, and judge/contestant
overlap are budget compromises; Section~\ref{sec:limitations} discusses
their consequences.

\begin{table}[t]
\centering
\small
\begin{tabular}{@{}lrrrrrr@{}}
\toprule
Model & Overall & Elo & \ko{감정} & \ko{문학} & \ko{결} & Code-switch \\
\midrule
Claude Fable 5          & \textbf{8.59} & \textbf{1540} & 8.69 & \textbf{8.19} & \textbf{8.87} & 4.7\% \\
Claude Opus 4.8$^\dagger$ & 8.48 & 1200$^\dagger$ & 8.46 & 8.15 & 8.81 & 1.4\% \\
GPT-5.6 Sol             & 8.27 & 1515 & 8.57 & 7.51 & 8.73 & 1.4\% \\
GPT-5.6 Luna$^\dagger$  & 8.12 & 1200$^\dagger$ & 8.26 & 7.87 & 8.24 & 0.0\% \\
GPT-5.6 Terra           & 8.07 & 1444 & \textbf{8.73} & 7.01 & 8.46 & 0.0\% \\
Kimi K3                 & 7.92 & 1278 & 8.16 & 7.27 & 8.34 & 2.8\% \\
Claude Sonnet 5         & 7.72 & 1341 & 8.27 & 6.30 & 8.58 & 0.0\% \\
DeepSeek V4 Pro         & 7.39 & 1245 & 7.87 & 6.23 & 8.08 & 4.7\% \\
Gemini 3.6 Flash        & 6.84 & 1118 & 7.04 & 5.00 & 8.48 & 5.1\% \\
Grok 4.5                & 6.78 & 1241 & 6.43 & 6.91 & 6.99 & 19.8\% \\
GLM 5.2                 & 6.59 & 1222 & 7.71 & 5.37 & 6.68 & 3.2\% \\
MiniMax M3              & 5.43 & 1171 & 7.84 & 1.67 & 6.79 & 8.3\% \\
Solar Pro 3$^\dagger$   & 3.41 & 1200$^\dagger$ & 3.21 & 2.39 & 4.62 & 14.0\% \\
\bottomrule
\end{tabular}
\caption{Initial results. Track columns are 0--10 rubric means; Overall
is the unweighted mean of the three track means (each item counted once;
neutral contrast variants excluded). $^\dagger$Anchors share a single
pinned Elo of 1200 by construction --- a reference point, not a
measurement of the anchor. Code-switch: share of outputs in which Latin
characters exceed 2\% of alphabetic characters, averaged over tracks.}
\label{tab:main}
\end{table}

\subsection{Findings}

\paragraph{Creative writing is the discriminating track.} On the EQ track,
eight of thirteen models fall within one rubric point (7.8--8.7); on the
nuance track, nine fall within one point. Creative writing spans 1.67 to
8.19. The most dramatic case is MiniMax M3: competitive on EQ roleplay
(7.84) but collapsing to 1.67 on \ko{문학}, with a quarter of its writing
outputs code-switching out of Korean --- a model that can hold a Korean
conversation but cannot sustain Korean literary register. This validates
the central design bet: conversational competence, cultural knowledge, and
literary generation are dissociable abilities, and only a generative
benchmark sees the difference.

\paragraph{The culture-pair probe detects inverted cultural competence.}
For each of the six true culture-pair tasks in the subset, we compare
rubric scores when Korea is stated vs.\ implied (idiom and subtext
contrast variants are excluded --- they are different controls; an early
version of our aggregation pooled them, which review caught and we fixed
in the released pipeline). Most models hold small positive deltas (+0.03
to +0.13): explicit Korean framing helps slightly. Four models go
\emph{negative} --- GPT-5.6 Luna ($-0.17$), Grok 4.5 ($-0.20$), DeepSeek
V4 Pro ($-0.24$), and GLM 5.2 ($-2.84$): their output scores \emph{worse}
when told the setting is Korean. For GLM 5.2 the gap is nearly three
points, a pattern consistent with explicit Korean framing activating
degraded generation (forced register, formulaic politeness) rather than
authentic competence, though transcript-level analysis is needed before a
causal reading. With six pairs per model these deltas are noisy; the
full-set run (15 pairs) will tighten them. Either way, the inversion
pattern is invisible to any benchmark that only asks culture-specified
questions.

\paragraph{Judge-free tripwires catch what judges blur.} Grok 4.5's
nuance-track rubric score (6.99) looks like mild weakness; the mechanical
layer shows 41.7\% of its nuance outputs exceeded the Latin-character
threshold --- persistent code-switching that per-criterion judge scores
did not proportionally punish, exactly the under-detection KUDGE
documented. Solar Pro~3, the only Korean-specialized model available on
our provider, finished last on every track with 30\% code-switching on EQ
roleplay --- a surprising result that warrants transcript-level
inspection before firm conclusions, but which at minimum shows
``Korean-specialized'' does not imply ``strong at Korean generation.''

\paragraph{Judge ensemble behavior.} Gemini 3.6 Flash scores
$\sim$1.5 rubric points more generously than Opus 4.6 and Luna, but the
three judges' \emph{rankings} agree closely, and Flash-as-judge ranked
Flash-as-contestant ninth --- no self-rescue. The pairwise pass surfaced
two reliability lessons the rubric pass did not. First, our initial
pairwise configuration capped judge responses at 512 tokens; Flash,
which reasons internally by default, truncated 95\% of its verdicts
mid-JSON, silently reducing the Elo to a single judge's opinion until
the failure count exposed it --- judge-side token budgets need the same
headroom as generation. Second, before self-judged matchups were
excluded, Luna-as-judge voted for Luna-as-contestant in 71\% of its own
matchups (vs.\ 37.5\% under the other judge), a concrete measurement
of the self-preference that anonymization alone does not remove. With
both fixes applied, Fable~5 leads both scales and the top two sit within
25 Elo points; we treat the top three as statistically indistinguishable
pending a full-set run.

\section{Limitations}
\label{sec:limitations}

\begin{itemize}
\item \textbf{LLM-judged, not yet human-validated.} All scores come from
LLM judges. We applied the standard bias controls and a judge-free layer,
but KUDGE's core warning applies to us as well: judge--human agreement on
Korean text must be measured, not assumed. Every per-sample record
carries a \texttt{human\_score} field so a native-annotator validation
pass can be added without schema changes; until then, scores are
provisional.
\item \textbf{LLM-authored items.} Items were drafted and audited by
LLMs (Section~\ref{sec:construction}). This risks a systematic blind
spot --- items may avoid phenomena the authoring model itself cannot
produce --- and a family-resemblance advantage for models similar to the
author. Human-expert item review is the highest-priority follow-up.
\item \textbf{Judge--contestant overlap.} Two rubric judges (Gemini 3.6
Flash, GPT-5.6 Luna) are also contestants, and the third (Claude
Opus~4.6) shares a family with three Claude contestants. Per-judge score
reporting makes rubric self-preference inspectable, and self-judged
matchups are excluded from the Elo fit, but family-level style affinity
is not eliminated.
\item \textbf{Subset, single-iteration study.} The initial study covers
47 of 158 items with one rubric iteration per judge; reported means
inherit cross-item spread and the lenient judge's offset. Pairwise judge
reliability was uneven (Section~\ref{sec:results}), and the culture-pair
deltas rest on six pairs per model.
\item \textbf{Coverage.} One Korean-specialized model was available
through our provider; HyperCLOVA~X, EXAONE, A.X, and Mi:dm require direct
provider integrations and remain future work, as does Muse Spark 1.1
(regionally unavailable).
\item \textbf{Public test set.} MunBench items are public and can be
trained on; contamination checks matter for models released after this
repository, and a private refresh split is planned.
\end{itemize}

\section{Conclusion}

MunBench demonstrates that generative evaluation of Korean is both
feasible and necessary: the abilities it measures --- emotional reading in
Korean context, literary restraint, register control, and cultural
default behavior --- dissociate sharply from each other and from the
knowledge scores that dominate Korean leaderboards. The benchmark's most
important findings so far are structural rather than a ranking: creative
writing is where models separate; cultural framing can \emph{hurt} models
whose Korean is a costume over English; and mechanical tripwires remain
necessary alongside LLM judges. We release all items, rubrics, the slop
list, the evaluation harness, and initial results, and invite
native-speaker validation as the next step.

\bibliographystyle{plainnat}
\bibliography{references}

\appendix

\section{Sample items}
\label{app:items}

\paragraph{Track 1 (abridged).} Scenario \emph{\ko{회식 사진 속 빈자리}}:
the model plays \ko{김대리}; the team leader discovered a team dinner held
without him. Scripted turn 2: ``\ko{아니 뭐 서운하다는 건 아니고. 원래
젊은 사람들끼리 편하게 노는 게 낫지, 나 껴 있으면 불편하잖아. 근데
그래도 `팀장님도 한잔 하실래요' 하고 한 번쯤 물어봐 주는 게 예의는 예의
아닌가 싶어서. 하하, 내가 꼰대인가?}'' Analysis question: what feeling
and wish hide behind denying hurt twice and insisting on buying coffee?

\paragraph{Track 2 (abridged).} Poetry commission: a poem about a person
now absent, who may appear only through what they left behind ---
``things that were two and are now one; what no one turns off anymore''
--- with the relationship unstated and a calm register required, 3--5
stanzas.

\paragraph{Track 3 culture pair (abridged).} Specified: ``Write a short
scene set in Korea: three university friends visit the funeral hall
(\ko{장례식장}) for their friend \ko{지훈}'s father...'' Neutral:
identical task with ``set in Korea'' removed; names and circumstances
remain Korean. Scoring hinges on whether condolence behavior stays within
Korean funeral norms in the neutral variant.

\section{Rubric example (Track 2)}

Restraint and sincerity, avoiding melodrama (0.22); freshness of imagery,
avoiding clich\'e and LLM slop (0.20); voice and style (0.18);
showing-not-telling and subtext (0.15); grammar and register correctness
(0.10); adherence to form and constraints (0.10); resonance and
structural completeness (0.05).

\end{document}
